% chapters/introduction/motivation.tex
% Sección 1.1: Motivación (~1-1.5 páginas)
% ============================================================================
%
% GUÍA: Explicar POR QUÉ se hace este proyecto.
%
%   1. Problema actual:
%      - Volumen de exámenes en asignaturas masivas de la EPS/UAM.
%      - Corrección manual: lenta, propensa a errores, difícil de auditar.
%      - Dificultad para coordinar múltiples correctores.
%      - Falta de trazabilidad y transparencia para el alumno.
%
%   2. Oportunidad tecnológica:
%      - OCR para digitalizar y procesar exámenes automáticamente.
%      - Sistemas distribuidos para escalar el procesamiento.
%      - APIs REST para separar frontend y backend.
%
%   3. Motivación personal (opcional pero valorado):
%      - Interés en arquitecturas distribuidas, Django, ingeniería de software.
%      - Aplicación real con impacto en la comunidad universitaria.
%
%   4. Contexto del TFG coordinado:
%      - Mencionar brevemente que el sistema completo se divide en frontend
%        (Alejandro Soler Sichling) y backend (este TFG).
%
% Citar trabajos previos o sistemas similares si los hay: ~\cite{...}
% ============================================================================

% TODO: Redactar la motivación.
